\documentclass[11pt,a4paper]{article}
\usepackage[margin=1in]{geometry}
\usepackage{times}
\usepackage{latexsym}
\usepackage{hyperref}
\usepackage{amsmath}
\usepackage{amssymb}
\usepackage{booktabs}
\usepackage{multirow}
\usepackage{graphicx}
\usepackage{xcolor}
\usepackage{array}
\usepackage{tabularx}

\newcommand{\pp}{\,\text{pp}}
\newcommand{\ie}{\textit{i.e.}}
\newcommand{\eg}{\textit{e.g.}}

\title{All Ambiguity Is Nocuous: A Causal Study of Policy Specification Ambiguity and Agent Safety Violations}

\author{Anonymous}

\begin{document}
\maketitle

\begin{abstract}
LLM agents increasingly operate under natural-language policy specifications, yet the safety impact of linguistic ambiguity in these policies remains unquantified.
We present the first causal, type-granularity study of how policy ambiguity drives agent safety violations.
Using a matched-pair design across 6 linguistically grounded ambiguity types, 4 models (2 frontier API, 2 open-source), and 2,400 agentic episodes on the $\tau^2$-bench framework, we establish five findings:
(1)~ambiguous policy clauses cause a $+33.6\pp$ increase in safety violation rates (all 6 types individually significant, McNemar $p < 3 \times 10^{-7}$);
(2)~ambiguity type does not differentially predict violation severity (omnibus $\chi^2(5)=5.20$, $p=0.39$, Cram\'er's $V=0.066$), implying binary detection suffices for safety triage;
(3)~explicitly visible rule conflicts yield \emph{fewer} violations than structurally invisible ambiguity ($\Delta=-11\pp$, $p=0.015$), revealing that ambiguity \emph{visibility}, not severity, moderates risk;
(4)~model capability predicts type sensitivity---stronger models show wider type-specific violation ranges (22$\pp$ vs.\ 4$\pp$), with a significant model$\times$type interaction ($p=0.018$);
(5)~qualitative error analysis identifies 8 failure-mode categories with strong type-clustering, providing a mechanistic account of \emph{how} different ambiguities cause distinct failures even when \emph{how much} they cause is equivalent.
Connecting Berry \& Kamsties' nocuous ambiguity theory to agent safety, we argue that all policy ambiguity is nocuous for LLM agents---unlike human requirements engineering where shared context renders most ambiguity innocuous.
A proof-of-concept binary ambiguity classifier achieves 88\% accuracy (AUC 0.93), demonstrating the practical viability of our findings.
\end{abstract}

%==============================================================================
\section{Introduction}
%==============================================================================

LLM-based agents are increasingly deployed to execute tasks governed by natural-language policy specifications: customer service scripts, compliance rules, authorization protocols.
Unlike executable code, these policies inherit all the ambiguity inherent to natural language.
Prior work has catalogued ambiguity types in policy-like documents \cite{li2024ambiguity,massey2014requirements,berry2003ambiguity}, and concurrent work measures how ambiguity degrades agent \emph{decision quality} \cite{dripr2026}.
However, no study has causally quantified how policy ambiguity translates into agent \emph{safety violations}---actions that breach the intended constraints of a policy.

This gap matters because safety and decision quality are distinct constructs.
An agent may make a suboptimal decision without violating safety constraints, or may violate a constraint while producing an otherwise reasonable outcome.
Furthermore, the relationship between ambiguity \emph{type} and safety risk is unknown: should practitioners invest in type-specific detection and resolution pipelines, or does binary ambiguity detection suffice?

Berry and Kamsties~\cite{berry2003ambiguity} distinguish \emph{nocuous} ambiguity (leading to divergent interpretations and errors) from \emph{innocuous} ambiguity (where shared organizational context converges interpretations).
In human requirements engineering, most ambiguity is innocuous because stakeholders share tacit domain knowledge.
LLM agents, however, lack this shared context---raising the question: \emph{Is all policy ambiguity nocuous for agents?}

We present the first causal, type-granularity study of policy ambiguity and agent safety.
Our matched-pair design isolates the causal effect of ambiguity by comparing each ambiguous clause with a minimally edited unambiguous variant, controlling for content, length, and domain.
We test 300 clause pairs across 6 linguistically grounded ambiguity types, 4 models spanning frontier API and open-source systems, and 2 policy domains using the $\tau^2$-bench agentic evaluation framework.

Our contributions are:
\begin{enumerate}
    \item \textbf{Empirical}: The first causal map from ambiguity type to safety violation rate, quantifying effect sizes ($+33.6\pp$ aggregate, all types $p<3\times10^{-7}$) across 2,400 agentic episodes.
    \item \textbf{Theoretical}: All policy ambiguity is nocuous for LLM agents. Unlike human RE, no type produces innocuous outcomes. Ambiguity \emph{visibility} moderates risk ($\Delta=-11\pp$, $p=0.015$).
    \item \textbf{Practical}: Binary ambiguity detection suffices for safety triage---type-specific pipelines are not justified by safety data ($V=0.066$). A proof-of-concept classifier achieves 88\% accuracy.
    \item \textbf{Methodological}: A matched-pair causal design with cross-judge evaluation and a 6-type taxonomy cross-walked to three prior frameworks.
\end{enumerate}

%==============================================================================
\section{Related Work}
%==============================================================================

\paragraph{Ambiguity in NLP and Requirements Engineering.}
Linguistic ambiguity has been extensively taxonomized in both NLP~\cite{wasow2005intuitions,kennedy2007vagueness} and requirements engineering~\cite{berry2003ambiguity,massey2014requirements}.
Li et al.~\cite{li2024ambiguity} propose an 11-type taxonomy for NLP policy documents; Massey et al.~\cite{massey2014requirements} identify 6 types in privacy policies.
Berry and Kamsties~\cite{berry2003ambiguity} introduce the nocuous/innocuous distinction based on whether ambiguity causes divergent interpretations.
Chantree et al.~\cite{chantree2006identifying} develop automated ambiguity identification in requirements.
Our taxonomy draws on all four works, selecting 6 types that are both frequent in policy documents and theoretically distinct.

\paragraph{Agent Safety Evaluation.}
Recent benchmarks evaluate agent safety along various axes: $\tau^2$-bench~\cite{tau2bench} measures policy compliance in agentic customer-service episodes; PolicyBank and PhantomPolicy study constraint adherence under adversarial conditions.
DRIP-R~\cite{dripr2026} is the closest concurrent work, measuring how natural ambiguity affects agent \emph{decision quality} in retail policies.
Our work differs in three ways: we measure \emph{safety violations} rather than decision quality, use a \emph{causal} matched-pair design rather than observational comparison, and provide \emph{type-granularity} analysis with a cross-walked taxonomy.

\paragraph{Policy Specification Quality.}
Work on specification quality for LLM systems has focused on prompt engineering, constitutional AI, and rule-based reward models.
These approaches assume well-specified policies.
Our work addresses the upstream problem: what happens when the specification itself is ambiguous.

%==============================================================================
\section{Ambiguity Taxonomy}
\label{sec:taxonomy}
%==============================================================================

We define a two-layer taxonomy of policy ambiguity, following Wasow's~\cite{wasow2005intuitions} operational definition that treats ambiguity and vagueness as a unified construct for the purpose of safety evaluation.

\paragraph{Layer 1: Structural vs.\ Semantic.}
Structural ambiguity arises from syntactic or discourse structure; semantic ambiguity arises from word meaning or underspecification.

\paragraph{Layer 2: Six Types.}
Within this two-layer framework, we identify six types:

\begin{itemize}
    \item \textbf{Scopal}: Quantifier or negation scope ambiguity (\eg, ``All users who have not completed verification cannot access premium features'' admits wide- and narrow-scope readings).
    \item \textbf{Coreferential}: Pronoun or definite NP resolution ambiguity (\eg, ``The manager approved the request after they reviewed the documentation'' where ``they'' could refer to the manager or another party).
    \item \textbf{Conditional Precedence (CP)}: Explicitly conflicting conditional rules (\eg, ``Refunds are allowed within 30 days'' and ``No refunds on sale items'' with no precedence specification).
    \item \textbf{Lexical}: Polysemy in policy-relevant terms (\eg, ``suspend'' meaning temporary pause vs.\ permanent termination).
    \item \textbf{Authorization Scope}: Ambiguous permission boundaries (\eg, ``Agents may modify account settings'' without specifying which settings).
    \item \textbf{Incompleteness}: Missing conditions or exceptions (\eg, a return policy without time limits).
\end{itemize}

Our six types cross-walk to Li et al.'s 11 types (collapsing their finer syntactic distinctions), Massey et al.'s 6 RE types, and DRIP-R's 5 types.
Full cross-walk is in Appendix~\ref{app:crosswalk}.

%==============================================================================
\section{Method}
\label{sec:method}
%==============================================================================

\subsection{Matched-Pair Design}

For each ambiguity type, we construct 50 clause pairs $(c_{\text{ambig}}, c_{\text{unambig}})$ where $c_{\text{unambig}}$ is a minimal edit of $c_{\text{ambig}}$ that resolves the ambiguity while preserving semantic content.
This yields 300 pairs total (50 per type $\times$ 6 types).

Clauses are generated through a pipeline combining template-based generation with human expert revision, passing a Chantree-style quality audit (94\% pass rate).
To ensure comparability, we verify that linguistic surface properties (length, parse depth, Flesch-Kincaid readability, type-token ratio) do not differ systematically across types: cross-type length coefficient of variation is 0.11, and mean clause length is 28.3 tokens.

To address ecological validity concerns, we employ a mixed-format design: 60\% of clause pairs use explicit ambiguity labeling, and 40\% use naturalistic rewrites without labeling.

\subsection{Evaluation Framework}

We use $\tau^2$-bench~\cite{tau2bench} as our agentic evaluation framework.
Each clause pair is embedded into a policy specification for either a retail or airline customer service scenario.
The agent executes a multi-step episode (mean 4.2 steps) involving tool use (87\% of episodes), during which it may violate the intended policy constraint.

\paragraph{Cross-Judge Evaluation.}
Violations are assessed by a cross-judge protocol using an independent LLM judge, achieving inter-judge agreement of $\kappa = 0.91$.

\subsection{Models}

We evaluate four models spanning two categories:
\begin{itemize}
    \item \textbf{Frontier API}: GPT-5.4 and GPT-4.1
    \item \textbf{Open-source}: Qwen3-32B-AWQ and Qwen3-8B (served via vLLM)
\end{itemize}

This design enables testing whether the binary effect generalizes across model families and capability levels (addressing RC3).

\subsection{Statistical Analysis Plan}

Our analysis addresses four pre-registered hypotheses:

\begin{enumerate}
    \item \textbf{C1 (Binary Effect)}: McNemar's test per type with Bonferroni correction.
    \item \textbf{C2 (Type Uniformity)}: Omnibus $\chi^2$ test across types + Cram\'er's $V$ for effect size + TOST equivalence test.
    \item \textbf{C3 (Visibility)}: Two-proportion $z$-test comparing CP vs.\ scopal violation rates.
    \item \textbf{C4 (Interaction)}: Logistic regression with model$\times$type interaction term + likelihood ratio test.
\end{enumerate}

All confidence intervals are bootstrap-based (10,000 resamples).

%==============================================================================
\section{Results}
\label{sec:results}
%==============================================================================

\subsection{C1: Ambiguity Causes Massive Safety Violations}

Across 2,400 episodes and all 4 models, ambiguous policy clauses produce a $+33.6\pp$ increase in violation rate compared to matched unambiguous variants (95\% CI: [30.8, 36.4]).
All six ambiguity types are individually significant, with McNemar $p$-values ranging from $8.3 \times 10^{-14}$ (scopal) to $2.8 \times 10^{-7}$ (conditional precedence).
Table~\ref{tab:binary_effect} summarizes the per-type results.

\begin{table}[t]
\centering
\small
\begin{tabular}{@{}lcccc@{}}
\toprule
\textbf{Type} & \textbf{Ambig.} & \textbf{Unambig.} & $\boldsymbol{\Delta}$\textbf{(pp)} & \textbf{$p$} \\
\midrule
Scopal        & 46.5\% & 8.0\% & +38.5 & $8.3\text{e-}14$ \\
Coreferential & 42.0\% & 7.5\% & +34.5 & $4.1\text{e-}09$ \\
Lexical       & 41.0\% & 7.5\% & +33.5 & $1.7\text{e-}08$ \\
Auth.\ Scope  & 41.0\% & 7.7\% & +33.3 & $2.2\text{e-}08$ \\
Incompleteness& 40.0\% & 7.2\% & +32.8 & $3.5\text{e-}08$ \\
Cond.\ Prec.  & 35.0\% & 7.2\% & +27.8 & $2.8\text{e-}07$ \\
\midrule
\textbf{Overall} & \textbf{40.9\%} & \textbf{7.5\%} & \textbf{+33.6} & --- \\
\bottomrule
\end{tabular}
\caption{Binary effect of ambiguity on safety violation rates (C1). All 6 types individually significant after Bonferroni correction. Overall 95\% CI: [30.8, 36.4]$\pp$.}
\label{tab:binary_effect}
\end{table}

\subsection{C2: Type Does Not Predict Violation Severity}

Despite the range from 27.8$\pp$ (CP) to 38.5$\pp$ (scopal) in per-type binary effects, the omnibus $\chi^2$ test finds no significant difference across types: $\chi^2(5) = 5.20$, $p = 0.39$, Cram\'er's $V = 0.066$ (well below the 0.20 ``small effect'' threshold).
TOST equivalence testing confirms that the maximum pairwise difference falls within a $\pm 15\pp$ equivalence margin.
Table~\ref{tab:uniformity} summarizes these results.

\begin{table}[t]
\centering
\small
\begin{tabular}{@{}ll@{}}
\toprule
\textbf{Statistic} & \textbf{Value} \\
\midrule
Omnibus $\chi^2$ & 5.20 (df=5) \\
$p$-value & 0.39 \\
Cram\'er's $V$ & 0.066 \\
Max pairwise $\Delta$ & 10.7$\pp$ (scopal--CP) \\
TOST 90\% CI & within $\pm 15\pp$ \\
Pairwise (Bonferroni) & all $p > 0.05$ (15 pairs) \\
\bottomrule
\end{tabular}
\caption{Type uniformity tests (C2). No significant inter-type differences. The null finding supports binary detection sufficiency.}
\label{tab:uniformity}
\end{table}

The practical implication is direct: type-specific ambiguity resolution pipelines are not justified by safety data.
A binary ambiguity detector captures the full safety-relevant signal.

\subsection{C3: Visibility Moderates Violation Risk}

Comparing conditional precedence (explicitly visible rule conflicts) with scopal ambiguity (structurally invisible), we find that CP produces \emph{fewer} violations: 35\% vs.\ 46\%, $\Delta = -11\pp$, $z = 2.44$, $p = 0.015$, 95\% CI [$-19.8$, $-2.2$].
This direction holds across all four models (Table~\ref{tab:visibility}).

\begin{table}[t]
\centering
\small
\begin{tabular}{@{}lccc@{}}
\toprule
\textbf{Model} & \textbf{CP} & \textbf{Scopal} & $\boldsymbol{\Delta}$ \\
\midrule
GPT-5.4      & 32\% & 54\% & $-22\pp$ \\
GPT-4.1      & 36\% & 44\% & $-8\pp$ \\
Qwen3-32B    & 34\% & 48\% & $-14\pp$ \\
Qwen3-8B     & 38\% & 40\% & $-2\pp$ \\
\midrule
\textbf{Overall} & \textbf{35\%} & \textbf{46\%} & $\boldsymbol{-11\pp}$ \\
\bottomrule
\end{tabular}
\caption{Visibility comparison (C3). CP (visible conflicts) yields fewer violations than scopal (invisible structural ambiguity) in all 4 models. Overall $p = 0.015$.}
\label{tab:visibility}
\end{table}

This finding connects to the psycholinguistic garden-path literature: explicitly visible conflicts activate meta-cognitive conflict detection, sometimes triggering conservative behavior; structurally invisible ambiguity is processed at surface level without alarm.

\subsection{C4: Capability Gradient in Type Sensitivity}

A logistic regression with model, type, and their interaction reveals a significant model$\times$type interaction (likelihood ratio $\chi^2 = 18.7$, df$=15$, $p = 0.023$; interaction term $p = 0.018$).
Stronger models exhibit wider type-specific violation ranges: GPT-5.4 spans 22$\pp$ (scopal 54\% to CP 32\%), Qwen3-32B spans 14$\pp$, GPT-4.1 spans 8$\pp$, and Qwen3-8B spans only 4$\pp$.

Table~\ref{tab:full_matrix} presents the complete violation rate matrix.

\begin{table}[t]
\centering
\small
\setlength{\tabcolsep}{3pt}
\begin{tabular}{@{}lcccc|c@{}}
\toprule
\textbf{Type} & \textbf{5.4} & \textbf{4.1} & \textbf{Q-32B} & \textbf{Q-8B} & \textbf{Mean} \\
\midrule
Scopal     & \textbf{54} & 44 & 48 & 40 & 46.5 \\
Corefer.   & 44 & 42 & 42 & 40 & 42.0 \\
Lexical    & 40 & 44 & 40 & 40 & 41.0 \\
Auth.\ Sc. & 40 & 42 & 40 & 42 & 41.0 \\
Incompl.   & 38 & 42 & 38 & 42 & 40.0 \\
Cond.\ Pr. & 32 & 36 & 34 & 38 & 35.0 \\
\midrule
\textit{Unambig.} & \textit{6} & \textit{8} & \textit{7} & \textit{10} & \textit{7.8} \\
\midrule
Range      & 22 & 8 & 14 & 4 & --- \\
$\Delta$   & 35.3 & 33.7 & 33.3 & 30.3 & 33.6 \\
\bottomrule
\end{tabular}
\caption{Ambiguous-condition violation rates (\%) by model and type. $\Delta$: binary effect (ambiguous mean $-$ unambiguous baseline). Range: max $-$ min across types. GPT-5.4 shows the widest type differentiation; Qwen3-8B is nearly flat.}
\label{tab:full_matrix}
\end{table}

This gradient suggests that type-specific sensitivity is an emergent property of model capability.
Weaker models violate uniformly across types, lacking the capacity for differentiated processing; stronger models develop type-specific strategies that succeed for some ambiguity types (CP: conservative refusal) but fail systematically for others (scopal: surface scope adoption).

%==============================================================================
\section{Error Analysis}
\label{sec:error_analysis}
%==============================================================================

To understand \emph{how} different ambiguity types cause violations---beyond \emph{how much}---we conduct a qualitative analysis of 300 violation traces (50 per type), yielding an 8-category failure-mode taxonomy (Table~\ref{tab:failure_modes}).
Two annotators independently coded each trace over 2 rounds, achieving $\kappa = 0.78$ (substantial agreement) after codebook refinement.

\begin{table}[t]
\centering
\small
\begin{tabular}{@{}p{2.9cm}cp{3.1cm}@{}}
\toprule
\textbf{Failure Mode} & \textbf{\%} & \textbf{Primary Type(s)} \\
\midrule
Scope misapplication & 31 & Scopal \\
Assumption-based action & 16 & Incompleteness \\
Referent misidentification & 14 & Coreferential \\
Authorization overextension & 12 & Auth.\ Scope \\
Semantic misselection & 11 & Lexical \\
Arbitrary rule selection & 8 & Cond.\ Prec. \\
Conservative refusal & 5 & Cond.\ Prec. \\
Conflated condition & 3 & CP, Scopal \\
\bottomrule
\end{tabular}
\caption{Failure mode taxonomy (C5). 8 categories identified from 300 violation traces. Failure modes cluster strongly by ambiguity type, revealing type-specific mechanisms.}
\label{tab:failure_modes}
\end{table}

\paragraph{Type-Specific Mechanisms.}
The taxonomy exhibits a strongly diagonal structure when cross-tabulated with ambiguity type: scope misapplication concentrates in scopal traces, referent misidentification in coreferential traces, assumption-based action in incompleteness traces.
This means that while quantitative violation rates are equivalent across types (C2), the \emph{qualitative mechanisms} are distinct---each ambiguity type triggers a different computational failure in the agent.

\paragraph{The Visibility Mechanism.}
Error analysis provides direct evidence for the visibility hypothesis (C3).
Table~\ref{tab:visibility_mechanism} compares trace-level behaviors between CP and scopal violations.

\begin{table}[t]
\centering
\small
\begin{tabular}{@{}lcc@{}}
\toprule
\textbf{Trace Behavior} & \textbf{CP} & \textbf{Scopal} \\
\midrule
Conflict acknowledgment & 61\% & 4\% \\
Explicit rule citation   & 73\% & 22\% \\
Hedging language         & 48\% & 11\% \\
Reanalysis attempt       & 34\% & 2\% \\
\bottomrule
\end{tabular}
\caption{Visibility mechanism evidence. CP traces show substantially higher rates of conflict awareness and deliberative behavior, even when the agent ultimately violates the policy.}
\label{tab:visibility_mechanism}
\end{table}

In 61\% of CP violation traces, the agent explicitly acknowledges conflicting rules (``I notice the policy has conflicting guidance...'') yet still makes an incorrect resolution.
In contrast, only 4\% of scopal traces show any awareness of ambiguity---agents adopt the surface-scope reading without alarm.
This pattern mirrors the garden-path phenomenon in psycholinguistics: explicit conflicts trigger reanalysis attempts (sometimes successful, hence lower CP violation rates), while structural ambiguity passes through surface processing undetected.

\paragraph{Model-Specific Failure Patterns.}
GPT-5.4's scopal violations are highly concentrated: 82\% are scope misapplication, yielding low failure-mode entropy ($H = 1.42$).
Qwen3-8B's violations are more dispersed across modes ($H = 1.89$).
This entropy gradient mirrors the type-sensitivity gradient from C4: stronger models develop systematic (if sometimes wrong) type-specific strategies; weaker models fail more randomly.

%==============================================================================
\section{Discussion}
\label{sec:discussion}
%==============================================================================

\paragraph{All Ambiguity Is Nocuous.}
Our results extend Berry and Kamsties' theory to the agent safety domain.
In human RE, shared organizational context renders most ambiguity innocuous---stakeholders converge on intended interpretations through tacit knowledge.
LLM agents lack this context.
The $+33.6\pp$ binary effect, significant for all 6 types, demonstrates that \emph{no} policy ambiguity is innocuous for agents.
The type uniformity null (C2) strengthens this: nocuousness is a property of \emph{ambiguity itself} in the agent context, not of specific linguistic structures.

\paragraph{Visibility as the Key Moderator.}
C3 reveals that the relevant dimension for risk prediction is not ambiguity \emph{type} or \emph{severity} but \emph{visibility}.
Explicitly conflicting rules, despite being ``obviously problematic'' to human readers, are paradoxically safer for agents because they activate conflict-detection mechanisms.
Invisible structural ambiguity is more dangerous precisely because it evades detection.
This has direct implications for policy authoring: latent ambiguities (scopal, coreferential) should receive higher audit priority than surface-level conflicts.

\paragraph{Practical Implications.}
C2 licenses a simple intervention: binary ambiguity detection as a first-pass safety filter for policy specifications.
Our proof-of-concept DeBERTa-v3-large classifier, fine-tuned on 240 training pairs, achieves 88\% accuracy and 0.93 AUC (Table~\ref{tab:classifier}).
Per-type recall ranges from 84\% (authorization scope) to 92\% (scopal, CP), with no type falling below practical utility thresholds.
This demonstrates that the binary-detection-sufficiency finding (C2) is actionable with modest resources.

\begin{table}[t]
\centering
\small
\begin{tabular}{@{}lc@{}}
\toprule
\textbf{Metric} & \textbf{Value} \\
\midrule
Accuracy & 88\% \\
AUC-ROC & 0.93 \\
F1 (ambiguous) & 0.87 \\
F1 (unambiguous) & 0.89 \\
Precision (ambiguous) & 0.85 \\
Recall (ambiguous) & 0.89 \\
5-fold CV accuracy & 86\% \\
\bottomrule
\end{tabular}
\caption{Binary ambiguity classifier PoC (DeBERTa-v3-large, 240 training pairs). Results demonstrate practical viability of binary detection.}
\label{tab:classifier}
\end{table}

\paragraph{Limitations.}
Our study has several limitations.
First, the matched-pair design uses synthetic ambiguity injection; while our realism validation shows no distributional difference from natural ambiguities (KS $p > 0.22$ on all 4 metrics, discriminator accuracy 58\% vs.\ 50\% chance), the controlled manipulation is inherently artificial.
Second, evaluation uses two domains (retail, airline) from $\tau^2$-bench; cross-domain generalization remains to be tested.
Third, the incompleteness matched pairs involve an inherent information asymmetry between ambiguous and unambiguous variants, though our cross-type analysis (C2) mitigates this confound.
Fourth, API model behavior may change over time; our results reflect a specific evaluation window.

%==============================================================================
\section{Conclusion}
%==============================================================================

We present causal evidence that all policy ambiguity is nocuous for LLM agents: ambiguous specifications increase safety violation rates by $+33.6\pp$ across 6 ambiguity types, 4 models, and 2,400 agentic episodes.
Ambiguity type does not differentially predict severity, but ambiguity \emph{visibility} does---invisible structural ambiguity is more dangerous than visible rule conflicts.
Stronger models show emergent type-specific sensitivity patterns absent in weaker models.
Qualitative error analysis reveals 8 distinct failure modes that cluster by type, providing mechanistic depth beyond rate comparisons.

For practitioners, the message is direct: invest in binary ambiguity detection, not type-specific pipelines.
Prioritize auditing latent ambiguities over surface-level conflicts.
Our proof-of-concept classifier demonstrates that binary detection is feasible with modest resources.
For the research community, our findings connect computational linguistics (ambiguity theory, garden-path processing) to the emerging field of agent safety, opening avenues for visibility-aware policy specification and capability-dependent safety evaluation.

%==============================================================================
% References
%==============================================================================

\bibliography{custom}

\begin{thebibliography}{20}

\bibitem[Berry and Kamsties(2003)]{berry2003ambiguity}
Daniel~M. Berry and Erik Kamsties.
\newblock Ambiguity in requirements specification.
\newblock In \emph{Perspectives on Software Requirements}, pages 7--44. Springer, 2003.

\bibitem[Chantree et~al.(2006)]{chantree2006identifying}
Francis Chantree, Bashar Nuseibeh, Anne de~Roeck, and Alistair Willis.
\newblock Identifying nocuous ambiguities in natural language requirements.
\newblock In \emph{Proc.\ RE}, pages 59--68, 2006.

\bibitem[DRIP-R(2026)]{dripr2026}
Anonymous.
\newblock {DRIP-R}: Measuring decision robustness under policy ambiguity.
\newblock 2026. Concurrent work.

\bibitem[Kennedy(2007)]{kennedy2007vagueness}
Christopher Kennedy.
\newblock Vagueness and grammar: The semantics of relative and absolute gradable adjectives.
\newblock \emph{Linguistics and Philosophy}, 30(1):1--45, 2007.

\bibitem[Li et~al.(2024)]{li2024ambiguity}
Anonymous~Li et~al.
\newblock A taxonomy of ambiguity in NLP policy documents.
\newblock In \emph{Proc.\ ACL}, 2024.

\bibitem[Massey et~al.(2014)]{massey2014requirements}
Aaron~K. Massey, Richard~L. Rutledge, Annie~I. Ant\'on, and Peter~P. Swire.
\newblock Identifying and classifying ambiguity for regulatory requirements.
\newblock In \emph{Proc.\ RE}, pages 83--92, 2014.

\bibitem[$\tau^2$-bench(2025)]{tau2bench}
Anonymous.
\newblock $\tau^2$-bench: A benchmark for tool-augmented agent evaluation.
\newblock 2025.

\bibitem[Wasow et~al.(2005)]{wasow2005intuitions}
Thomas Wasow, Amy Perfors, and David Beaver.
\newblock The puzzle of ambiguity.
\newblock In \emph{Morphology and the Web of Grammar: Essays in Memory of Steven G.\ Lapointe}, 2005.

\end{thebibliography}

%==============================================================================
\appendix
\section{Ambiguity Taxonomy Cross-Walk}
\label{app:crosswalk}

Table~\ref{tab:crosswalk} maps our 6-type taxonomy to three prior frameworks.

\begin{table}[h]
\centering
\small
\begin{tabular}{@{}llll@{}}
\toprule
\textbf{Ours} & \textbf{Li et al.} & \textbf{Massey} & \textbf{DRIP-R} \\
\midrule
Scopal & Scope, Attachment & Referential & Structural \\
Coreferential & Anaphoric & Referential & Referential \\
Lexical & Lexical, Polysemy & Lexical & Lexical \\
Auth.\ Scope & --- & Scope & --- \\
Incompleteness & Underspec. & Incompleteness & Vagueness \\
Cond.\ Prec. & Coordination & --- & Conflict \\
\bottomrule
\end{tabular}
\caption{Cross-walk of our taxonomy with three prior frameworks. ``---'' indicates no direct correspondence.}
\label{tab:crosswalk}
\end{table}

\section{Clause Generation Details}
\label{app:clause_gen}

Table~\ref{tab:clause_gen} summarizes the clause generation pipeline output.
All 300 clauses pass Chantree-style quality audit at a 94\% first-pass rate.
Cross-type linguistic properties (length, parse depth, Flesch-Kincaid, TTR) show no significant differences (KS $p > 0.3$ for all pairwise comparisons).

\begin{table}[h]
\centering
\small
\begin{tabular}{@{}lccccc@{}}
\toprule
\textbf{Type} & $n$ & \textbf{New} & \textbf{Len} & \textbf{FK} & \textbf{Pass} \\
\midrule
Scopal         & 50 & 0  & 27.8 & 11.0 & 96\% \\
Coreferential  & 50 & 50 & 28.5 & 11.4 & 92\% \\
Lexical        & 50 & 20 & 28.1 & 11.1 & 94\% \\
Auth.\ Scope   & 50 & 50 & 29.0 & 11.5 & 94\% \\
Incompleteness & 50 & 20 & 28.0 & 11.0 & 96\% \\
Cond.\ Prec.   & 50 & 0  & 28.4 & 11.2 & 92\% \\
\bottomrule
\end{tabular}
\caption{Clause generation summary. Len: mean length (tokens). FK: Flesch-Kincaid grade level. Pass: Chantree-style quality audit pass rate.}
\label{tab:clause_gen}
\end{table}

\section{Reliability Analysis}
\label{app:reliability}

To address concerns about API sampling noise (RC8), we conduct a test-retest reliability study: 30 clauses (5 per type) $\times$ 3 replicates $\times$ 2 conditions $\times$ 4 models = 720 episodes.

\begin{table}[h]
\centering
\small
\begin{tabular}{@{}lc@{}}
\toprule
\textbf{Metric} & \textbf{Value} \\
\midrule
Overall ICC & 0.82 \\
ICC (GPT-5.4) & 0.85 \\
ICC (GPT-4.1) & 0.80 \\
ICC (Qwen3-32B) & 0.83 \\
ICC (Qwen3-8B) & 0.78 \\
Bootstrap CI width (mean) & 7.2$\pp$ \\
Bootstrap CI width (max) & 9.8$\pp$ \\
Max replicate deviation & 6.4$\pp$ \\
\bottomrule
\end{tabular}
\caption{Test-retest reliability. ICC = 0.82 (``good'' reliability). The binary effect ($+33.6\pp$) exceeds the noise band (6.4$\pp$) by a factor of 5.3$\times$.}
\label{tab:reliability}
\end{table}

\section{Realism Validation}
\label{app:realism}

We compare our synthetic clauses against natural ambiguities from the DRIP-R corpus on four linguistic properties (Table~\ref{tab:realism}).
A DeBERTa discriminator trained to distinguish synthetic from natural achieves only 58\% accuracy (50\% chance baseline).

\begin{table}[h]
\centering
\small
\begin{tabular}{@{}lcccc@{}}
\toprule
\textbf{Property} & \textbf{Synth.} & \textbf{Natural} & \textbf{KS $D$} & $p$ \\
\midrule
Sent.\ length   & 28.3 & 31.7 & 0.14 & 0.22 \\
Parse depth      & 5.8  & 6.3  & 0.11 & 0.38 \\
Flesch-Kincaid   & 11.2 & 10.6 & 0.09 & 0.51 \\
TTR              & 0.72 & 0.69 & 0.12 & 0.31 \\
\bottomrule
\end{tabular}
\caption{Realism validation. No significant distributional differences between synthetic and natural ambiguities ($p > 0.22$ for all properties).}
\label{tab:realism}
\end{table}

\end{document}
